\chapter{Evaluation Methodology}

% =====================================================================
% Define QUÉ se compara y CÓMO se mide. Ningún resultado numérico en este
% capítulo. Toda métrica presentada en el Capítulo 6 debe estar definida aquí,
% y ninguna definición de métrica debe aparecer allí.
% =====================================================================

\section{Evaluated Configurations and Baselines}
\emph{Definir las cuatro configuraciones comparadas: referencia de acceso celular
exclusivo, referencia de acceso inalámbrico exclusivo, sistema híbrido con la
política predictiva, y sistema híbrido con la política predictiva y la extensión
por plazos. Indicar qué representa cada una como punto del espacio de compromiso
y qué variable única la separa de su referencia.}

\section{Simulation Scenarios and Operating Conditions}
\emph{Definir las condiciones bajo las que se evalúan las configuraciones:
escenario de movilidad, duración, niveles de demanda de tráfico de fondo y
condiciones de propagación consideradas. Declarar qué se mantiene constante entre
todas las ejecuciones.}

\section{Delivery and Reliability Metrics}
\emph{Definir formalmente las métricas de entrega: proporción de datos entregados,
pérdidas por descarte en memoria frente a pérdidas en el canal, y reparto de la
carga entre interfaces. Precisar el criterio con que se considera entregado un
dato.}

\section{Latency and Information Freshness Metrics}
\emph{Definir el retardo extremo a extremo y, de forma separada, la frescura del
dato en el momento de su recepción, junto con la proporción de datos que llegan
dentro de su plazo de validez. Es la métrica que da sentido a la extensión por
plazos y debe quedar definida sin ambigüedad.}

\section{Energy Metrics}
\emph{Definir las magnitudes energéticas comparadas: energía total consumida por
dominio, potencia media y energía por unidad de información útil entregada. Esta
última es la que permite comparar configuraciones que no entregan el mismo
volumen de datos.}

\section{Experimental Design and Statistical Treatment}
\emph{Documentar el diseño experimental: número de repeticiones por
configuración, gestión de semillas, tratamiento del transitorio inicial, y
estadísticos e intervalos empleados para presentar los resultados. Declarar el
criterio con el que se considera significativa una diferencia entre
configuraciones.}
